\documentclass[11pt]{article}
\usepackage[a4paper,margin=1in]{geometry}
\usepackage{amsmath,amssymb,amsthm,mathtools,stmaryrd}
\usepackage{hyperref,booktabs,longtable,array}
\hypersetup{colorlinks=true,urlcolor=blue,linkcolor=blue}
\newtheorem{proposition}{Proposition}
\newtheorem{definition}{Definition}
\newtheorem{observation}{Observation}
\newcommand{\Str}{\mathsf{Str}}
\newcommand{\Play}{\mathsf{play}}
\newcommand{\End}{\mathsf{End}}
\newcommand{\Q}{\mathsf{Q}}
\title{A role-symmetric logic and dependent calculus:\\research assessment of K-FOG}
\author{Research notes for the KFOL workspace}
\date{28 October 2026}
\begin{document}
\maketitle

\noindent\textbf{Assessment.}
The central object is mathematically viable, and there is a particularly direct
precedent: Abramsky's $L_A$. The role-indexed quantifier idea, including independent
per-agent valuations, should not be claimed as new. There are also existing
multirole proof calculi and role-indexed dependent session quantifiers. These
findings support the idea's coherence, while substantially changing its novelty
position. I have not established that the desired full dependent proof calculus
already exists, nor established a novel theorem for it.

The object of interest here is intrinsically multi-role interaction, rather than
a language organized around individual forceability. The latter is one optional
observation of a game and is not used to define the proposed core. Several
forceability claims are audited in the appendix because the manuscript makes them.
The original manuscript has not been edited.

\section{Which existing objects actually match?}

\paragraph{The closest semantic match: Abramsky's $L_A$.}
\href{https://www.cs.ox.ac.uk/people/samson.abramsky/admproc.pdf}{A Compositional Game
Semantics for Multi-Agent Logics of Partial Information} (2007), especially
Sections 2.4, 3.3--3.4, and 4.4--4.6, gives role-owned quantifiers, role
permutations, strategy interaction, per-agent valuations, and partial information.
It treats formulas as interaction structures. Its Section 6 leaves proof theory
open; that is a statement about that paper, not a claim that the problem remains
open today.

For the manuscript's broadcast prenex fragment the comparison can be made exact,
using its intended full-history strategies rather than the inconsistent
free-variable wording in its strategy definition.
Translate each $P_i x$ to a choice particle owned by $i$. Package the terminal
$\langle\varphi_0,\ldots,\varphi_{K-1}\rangle$ as an interpreted terminal atom
$a_{\vec\varphi}(\vec x)$ whose $i$-valuation is
\[
  v_i(a_{\vec\varphi},\mathcal M,s)
  :=[\mathcal M,s\models\varphi_i].
\]
This is a definitional interpretation of a new terminal symbol, not a claim that
every tuple is already an atom in the original first-order signature.
It gives the candidate correspondence
\[
 v_i\bigl(\Play(T\Phi,T\sigma)\bigr)
       = \llbracket\Phi\rrbracket_{\mathcal M,s,\sigma,i}.
 \tag{1}
\]
For a fixed broadcast prefix, strategies in both presentations can be restricted
to the same history-to-choice functions. An induction on the prefix then proves
(1): the same owner selects the same value, and the terminal clause is true by
definition. This is a correspondence of full evaluations, before choosing a
success predicate. For arbitrary information annotations, one must additionally
prove that the two presentations admit exactly the same observation-uniform
strategies. The manuscript's inconsistent slash conventions cannot be translated
as a single convention without repair.

\paragraph{A close algebraic antecedent.}
\href{https://pure.uva.nl/ws/files/1012623/156995_334017.pdf}{Olde Loohuis and Venema,
Logics and Algebras for Multiple Players} (2010) uses player-owned connectives,
role switches, and valuations declaring arbitrary subsets of players winners.
It develops equational proof systems and representation results. Its chosen
equivalence observes winning-strategy sets; retaining the complete interaction
object requires a finer equivalence. This is a direct antecedent for the
multi-outcome idea, even though it is not the desired dependent term calculus.

\paragraph{A close proof-theoretic antecedent.}
\href{https://arxiv.org/pdf/2309.01317}{Xi and Wu, Multirole Logic and Multiparty
Channels}, with earlier work in
\href{https://arxiv.org/abs/1703.06391}{2017} and
\href{https://arxiv.org/abs/1611.08888}{2016}, develops role-indexed logical
operations and multiparty cut elimination. Its first-order quantifiers are indexed
by ultrafilters; on a finite role set these are principal. It includes classical,
intuitionistic, and linear variants. This is evidence for a real proof theory of
role-indexed logic, not merely a protocol analogy. It does not automatically
validate the manuscript's unrestricted payoff tuples or proposed compositions.

\paragraph{The dependent quantifier already has an implementation precedent.}
\href{https://arxiv.org/pdf/1808.00077}{Wu and Xi, Multiparty Dependent Session
Types} (2018), Figure 7 and Section 4, has
$\mathsf{quan}(r,\lambda a:\sigma.\pi)$. Its interpretation varies by endpoint,
with the designated role interpreted universally and the others existentially
in that presentation. The polarity convention differs from an owner-as-witness
presentation; the neutral role-indexed mechanism is already there. The calculus
separates statics and dynamics, so this result alone is not an identification with
a full CIC-style dependent calculus.

\paragraph{Other comparisons a type-theory claim must survive.}
There are existing game models of dependent type theory, including
\href{https://arxiv.org/abs/1610.01669}{Yamada's work} and
\href{https://doi.org/10.1016/j.ic.2018.02.015}{Abramsky--Jagadeesan--V\'ak\'ar}.
Those results show that dependent game semantics itself is not new. They do not
by themselves supply this particular multi-role object language.
\href{https://arxiv.org/abs/2212.07735}{Escard\'o and Oliva} model dependent game
trees and mechanize strategy constructions in Agda. Using dependent types to
describe games is therefore also an existing direction.
\href{https://arxiv.org/abs/1603.04641}{Open games} supply a compositional treatment
of game-theoretic interaction and equilibria; a composition claim should compare
the preserved behavior, not merely the use of sequential and parallel notation.

For the dependent information discipline, compare
\href{https://arxiv.org/abs/2201.11040}{A Dependent Dependency Calculus} and
\href{https://arxiv.org/abs/2011.15021}{Multimodal Dependent Type Theory}. The former
integrates dependency analysis into dependent typing; the latter gives a general
framework for interacting modalities with consistency and canonicity results.
Neither is automatically an identification with neutral multi-party choice.
The current \href{https://github.com/kami-language/kami-theory}{Kami theory
repository} explicitly describes a distributed language based on MTT with full
dependent types, so a broad claim about being the first dependent multi-location
calculus would be particularly unsafe.

\section{The object to retain}

A useful semantic object is
\[
 G=\bigl(\text{typed interaction structure},
          (\Str_i(G))_{i<K},\Play_G,(B_i)_{i<K}\bigr).
 \tag{2}
\]
Here $\Play_G$ resolves a profile of strategies into a complete history, and $B_i$
is the $i$-th terminal meaning on that history. In the manuscript $B_i$ is
Boolean-valued. For a dependent theory it can instead be a family of propositions
or evidence types. General objectives remain general: no one-hot or common-goal
condition is imposed on (2).

No strategy is selected or declared successful in the definition of (2).
One can subsequently ask about proofs, partial implementations, refinement,
equilibria, or forceability. Keeping these separate prevents the latter from
silently becoming the whole meaning of a formula.

For a fixed prefix, number occurrences $m=1,\ldots,n$ and attach an owner $r_m$,
an action type, and an observation map
\[
 o_m:H_m\longrightarrow O_m.
\]
For fixed nondependent action types $A_m$, pure strategies are exactly
\[
 \Str_i(G)=\prod_{m:r_m=i}(O_m\longrightarrow A_m).
 \tag{3}
\]
This gives a precise meaning to a single strategy for a role: a family of local
policies, one per occurrence. Additional coherence requirements must be stated
as additional requirements; packaging the policies into one function does not
create them.

In a dependent version, start with action types that descend to observations:
there is a family $\bar A_m:O_m\to\mathcal U$ and
\[
 A_m(h)=\bar A_m(o_m h),\qquad
 f_m:\prod_{o:O_m}\bar A_m(o).
 \tag{4}
\]
This avoids assigning a player a choice type that varies with an unavailable
secret. A richer theory could allow abstract hidden witnesses, but would need
explicit packaging or a modal discipline rather than an unstated coercion.

When observations are sets of earlier columns $K_m$, perfect recall has a simple
checkable sufficient condition: whenever $b<m$ have the same owner,
\[
 K_b\cup\{b\}\subseteq K_m.
 \tag{5}
\]
The unrestricted occurrence-local calculus and the perfect-recall calculus are
both legitimate objects. They are different objects. The manuscript currently
uses properties of both without fixing which one is intended.

\section{Three precise foundational facts}

\begin{proposition}[Evaluation and change of schedule]
For an acyclic finite dependency structure, total deterministic local policies
of the form (3) determine a unique complete history. Changing to another
topological order of the same dependency structure does not change that history.
\end{proposition}
\begin{proof}
For any event, its predecessors determine its observation and hence its action.
Induct over the dependency order. Each event receives the same inputs in every
topological execution, so it produces the same value. This argument uses the same
event identities, action types, and policies; merely rearranging a broadcast
prefix without preserving observations does not meet the hypothesis.
\end{proof}

\begin{proposition}[Role equivariance]
A permutation of roles transports (2), its strategies, and its complete
evaluations equivariantly.
\end{proposition}
\begin{proof}
Relabel owners, reorder the terminal family, and transport each occurrence's
policy along the corresponding occurrence bijection. The values selected at
corresponding events are equal by induction. Thus the resulting terminal vector
is permuted. If policies take syntax as an argument, the transport must also
rename that argument; reordering the profile coordinates alone is insufficient.
\end{proof}

\begin{proposition}[No single-role permutation is general duality]
For $K>2$, no permutation of the $K$ individual roles implements Boolean
complement on every one-hot terminal outcome.
\end{proposition}
\begin{proof}
A permutation takes a one-hot vector to a one-hot vector. Its Boolean complement
has $K-1$ nonzero coordinates. These cannot be equal when $K>2$.
\end{proof}

The last fact explains a structural distinction. For local producer/consumer
behavior, if an endpoint owns a set $R$ of roles, the complementary behavior at
every role-owned choice belongs to the complementary set $\bar R$. The complement
of a singleton is not a singleton when $K>2$. Multiparty interaction therefore
needs a coherent family of endpoints, often indexed by a partition of roles.
This does \emph{not} justify defining logical negation by complementing role sets:
the interaction laws, identity, and cut must still be checked. The multirole
literature itself identifies failures of naive negation under multiparty cut.

\section{A small dependent kernel worth making exact}

Take a terminating dependent data calculus as a base. Initially all communicated
data and their types are public. Use a neutral global constructor
\[
  \Q_r(x:A).G
\]
with the continuation $G$ allowed to depend on the chosen value. Role labels
remain part of the object language. A local interpretation at role $i$ can be
defined, using the owner-as-witness convention, by
\[
 L_i(\Q_r(x:A).G)=
 \begin{cases}
   \displaystyle\sum_{x:A} L_i(G),&i=r,\\[3pt]
   \displaystyle\prod_{x:A} L_i(G),&i\ne r.
 \end{cases}
 \tag{6}
\]
This is a concrete constructive formulation of local witness production and
local input consumption. It is not a claim of novelty: it is close to the
dependent session quantifier mechanism already discussed. For ordinary behavior
types set $L_i(\End)=\mathbf1$. For certified endpoint types, optionally use an
evidence family $L_i(\End)=B_i$. The global object and its ordinary behaviors
exist independently of whether any certified endpoints exist.

For example, with inhabited $X$ and positive natural numbers $\mathbb N_{>0}$,
\[
 \Q_A(n:\mathbb N_{>0}).\;
 \Q_B(v:\mathsf{Vec}(X,n)).\;
 \Q_C(j:\mathsf{Fin}(n)).\;
 \End(B_A(n,v,j),B_B(n,v,j),B_C(n,v,j))
 \tag{7}
\]
contains an actual dependency of later choice types on earlier choices. At $A$
the first constructor is a dependent pair; at $B$ and $C$ it is a dependent
function. Subsequent constructors make different roles the producer. The tuple
of terminal families is retained globally.

The interaction reduction at the head choice is explicit:
\[
 \mathsf{run}\bigl(\langle a,t_r\rangle,
                     (\lambda x.t_i)_{i\ne r}\bigr)
 \longrightarrow
 \mathsf{run}\bigl(t_r,(t_i[a/x])_{i\ne r}\bigr).
 \tag{8}
\]
More generally apply arbitrary dependent-function terms to $a$; lambdas merely
display the beta step. The chosen value is substituted in all continuation
types as well as their terms.

\begin{proposition}[Dependent interaction for a fixed finite telescope]
Given local inhabitants of all the types (6), interaction produces one common
history. If terminal types are $B_i$, it also produces an inhabitant of each
$B_i$ at that same history. There is one head communication per global binder.
\end{proposition}
\begin{proof}
At a head owned by $r$, eliminate the dependent pair at endpoint $r$, obtaining
$a:A$ and its continuation. Apply every other endpoint's dependent function to
$a$. The results have exactly the local types of the continuation with $a$
substituted. Recurse on the shorter telescope. At $\End$ the endpoint terms
already inhabit the terminal families. Termination of the communications follows
from the decreasing telescope length; normalization of arbitrary data-level
computations is a separate property of the base calculus.
\end{proof}

This is already stronger than a Gallina-looking let-chain as an account of
typed computation. But it does not yet give implication between interacting
objects, higher-order interaction, a proof theory for all of (2), or a new
foundation for dependent mathematics. Those are separate claims.

\section{What would be a substantial research contribution?}

The most useful candidate is a \emph{dependent multiparty substitution/cut
theorem for the full interaction object}. It would give typed open terms,
interfaces, identity, application/composition, and a reduction semantics whose
meaning preserves role ownership and dependency. It should retain arbitrary
terminal families and the behavior of each role, rather than prove a theorem
only about one Boolean projection.

Here are concrete proof obligations, in an order that avoids claiming more than
has been built:
\begin{enumerate}
\item \textbf{Correspondence to the existing semantic object.} Formalize (1),
      including the admissible strategy families. This settles exactly which
      part of the manuscript is an existing fragment and which is an extension.
\item \textbf{An actual dependent term calculus.} Specify formation, introduction,
      elimination, substitution, and judgmental equality for a small neutral
      choice fragment. Prove (8) preserves typing and the appropriate observations.
      Merely defining a datatype of games inside CIC does not establish a new
      object type theory.
\item \textbf{Identity and implication.} Define what a term using $G$ to implement
      $H$ is, with its role-indexed interface. The manuscript's material implication
      between terminal goals is not such a definition: it supplies no application
      or substitution operation. State the required coherence at shared terminals
      instead of importing an identity axiom for arbitrary goal tuples.
\item \textbf{Composition of open terms.} Prove substitution is associative up to
      the declared equality and gives identity laws. For an information-aware
      extension, require an interface substitution to respect observation
      equivalences. In elementary terms, if $h\sim_i h'$, then the substituted
      input views available to $i$ must still agree. This condition composes and
      preserves uniform local policies. Extending it to dependent types is a
      substantive part of the design.
\item \textbf{Computation and checking.} Prove normalization or an appropriate
      weaker result for the chosen fragment, and establish decidable checking
      under explicit assumptions on the base calculus. Cut elimination alone
      does not imply confluence. Classical proof reduction has well-known
      confluence difficulties; see
      \href{https://drops.dagstuhl.de/entities/document/10.4230/LIPIcs.CSL.2012.320}{Hetzl
      and Stra\ss burger}. Ownership may help orient interactions, but is not a
      proof that every higher-order critical pair is joinable.
\item \textbf{A precise conservativity claim.} Give a translation of an ordinary
      dependent fragment and prove the claimed preservation and reflection.
      Greater expressiveness about interaction does not automatically mean
      greater proof-theoretic strength or a stronger consistency assumption.
\end{enumerate}

The first two items are reliable starting points and are not, on their own,
compelling novelty claims. The remaining items are candidate research results.
I have not proved them for a specified new calculus or established that no
existing system proves an equivalent theorem. Comparison must be against the
multirole systems, dependent session systems, and modal/dependency calculi above.

One subtlety in the existing intuitionistic multirole presentation is that its
intuitionistic judgment is parameterized by a fixed ultrafilter. Over finite
roles that parameter selects a principal role. An equivariant family of such
judgments is different from one fixed judgment with a privileged role. Whether
the desired theory needs a single fully neutral judgment, or an explicitly
role-indexed family of constructive judgments, is an exact design issue. Neither
choice should be concealed by notation.

I would initially keep universes predicative and interaction finite, with
dependent data and a small genuine composition fragment. General recursion,
impredicativity, full inductive families of interactive types, and arbitrary
information disclosure can be separate extensions. A small proof of a meaningful
cut theorem would settle more than a large collection of encodings or
complexity reductions.

\appendix
\section{Manuscript audit}

The following audit distinguishes false claims from repairable definitions and
results that are essentially correct under explicit hypotheses. File locations
refer to the unchanged manuscript in this workspace.

\subsection{Claims that survive}

\begin{itemize}
\item Evaluation by recursion on a finite prefix is well-defined for total
      strategies and nonempty action domains. If the domain is infinite, the
      game is finite-horizon, not necessarily a finite game tree.
\item Role equivariance is correct after transporting syntax-indexed strategies.
\item The two-player, complementary-payoff prenex embedding has the expected
      projection and role-swap laws in the broadcast fragment, after fixing the
      seat-indexing error described below.
\item The broadcast individual-projection decision problem is PSPACE-complete
      already over a fixed Boolean structure. This is a model-checking result,
      not a decidability claim about validity over all structures.
\item Joining total functional strategy graphs on an acyclic dependency structure
      computes the unique play. This retains counterfactual strategies only if
      the full graphs are retained, not just their resulting row.
\item The tensor/leaky product forceability equality for disjoint player blocks,
      unchanged component objectives, total pure strategies, and broadcast
      component games is repairable. To extract a strategy on the right, freeze
      one whole left-hand history consistently. The proof's claim that arbitrary
      history-by-history extensions can be chosen is too strong. Its first
      direction also needs attention to the enlarged opponent strategy sets.
\end{itemize}

These are useful facts but do not make individual forceability the intended core.

\subsection{The abstraction theorem is incorrectly indexed and overclaims uniqueness}

In \texttt{2FOL.tex:29}, $\mathsf A_i(\theta)$ always has terminal
$\langle\theta,\neg\theta\rangle$. For $i=1$ its $i$-projection is
$\neg\theta$, contradicting the stated retract. Put $\theta$ in seat $i$ and
$\neg\theta$ in the other seat.

The broad uniqueness assertion in \texttt{appendix-proofs.tex} does not follow
from projection laws. Semantically equivalent redundant structure can be
inserted. Moreover, the appendix uses Boolean game constructors absent from the
core grammar. Pointwise conjunction of complementary terminal pairs gives
$\langle a\land b,\neg a\land\neg b\rangle$, whose second component is not
$\neg(a\land b)$. The proof assumes the connective compatibility it needs to
establish. A uniqueness statement restricted to a stipulated constructor-by-
constructor translation is a different, much narrower statement.

\subsection{Three incompatible information conventions}

\texttt{info.tex} starts with roles hidden from a newly produced value, then
defines independence from earlier variables at the current choice.
\texttt{session.tex} uses the outgoing recipient convention;
\texttt{complexity.tex} uses incoming variable independence. These are distinct
annotations. Fix one observation-map semantics and translate any surface sugar
into it explicitly.

\texttt{KFOL.tex:57} describes strategies using assignments only to the free
variables of the remaining formula. That can omit earlier observations which
are unused by terminal predicates but usable as signals. Use complete histories
and then project to the declared observations.

\subsection{The DQBF reduction is false as written}

The reduction in \texttt{complexity.tex:55} hides only the disallowed universal
variables. Earlier existential choices remain visible and can carry forbidden
information. Consider the false DQBF
\[
 \forall y;\quad x_1(y),\ x_2();\quad
 (x_1=y)\land(x_2=y).
\]
The manuscript's encoding permits $x_1:=y$ followed by $x_2:=x_1$, so the encoded
player wins. The intended $x_2$ must be constant and cannot always equal $y$.
For occurrence-local information, repair the hidden set to
\[
 H_j=(\{y_1,\ldots,y_n\}\setminus D_j)
            \cup\{x_1,\ldots,x_{j-1}\}.
 \tag{9}
\]
This repair permits forgetting. It does not prove the theorem under strict
single-player perfect recall.

For the corrected Boolean finite-prefix calculus with arbitrary occurrence-local
views, individual forceability is in NEXPTIME as well as hard for it: guess the
exponential-size local truth tables and enumerate all opponent action sequences.
A losing sequence is realizable by opponent policies constant at each occurrence,
so universal opponent-strategy tables need not themselves be enumerated.
This upper-bound argument assumes unrestricted nonempty Boolean moves and a
polynomial-size Boolean terminal formula.

With perfect recall for the one forcing role and fixed incoming observation
masks, its observed opponent variables form nested sets. Earlier own choices
can be inlined as functions of the remembered observations. Place each universal
input immediately before the first forcing move that observes it, and place
never-observed inputs at the end. This yields an ordinary QBF of polynomial size,
giving a PSPACE bound for that restricted problem. This is compatible with the
known collapse of perfect-recall IF truth conditions to first-order truth
conditions; see the primary
\href{https://api.pageplace.de/preview/DT0400.9781139065382_A23865990/preview-9781139065382_A23865990.pdf}{Mann--Sandu--Sevenster
book preview}, discussion of Chapter 6. It does not equate full games, all
players' perspectives, or distributed coalitions with a single scalar formula.

\subsection{Nash and individual rationality}

\texttt{solutions.tex:134} says IR-NE may fail to exist when pure Nash equilibria
exist and every player can force a win. Under the given definitions:
\[
 \mathsf{NE}(\sigma)\quad\Longrightarrow\quad\mathsf{IR}(\sigma).
\]
If $i$ could force 1 but received 0, switching to that forcing strategy would
be a profitable unilateral deviation. Hence IR-NE and NE are the same set here.
Further, if every $i$ has a forcing strategy, choose one for each role. Their
joint profile wins for everyone, because each guarantee applies against all
others. That profile is automatically Nash for Boolean payoffs. This contradicts
the accompanying discussion of independently forcing profiles.

The definitions of pure NE and weak dominance themselves are standard and usable.
For finite perfect-information games with nonempty finite action sets, pure
subgame-perfect equilibria already exist by backward induction; hidden-information
extensions can lack pure NE.

\subsection{Equal objectives and negation}

Identical terminal objectives do not turn individual adversarial quantification
into collective existence. Over the Boolean domain,
\[
 P_0x.\langle x=0,x=0\rangle
\]
has individual forceability $(1,0)$ and is jointly realizable. Complementing both
goals still gives individual forceability $(1,0)$, by choosing $x=1$.
Thus the cooperative-negation remark in \texttt{KFOL.tex:253} is false.

The general no-De-Morgan conclusion is correct, but its displayed recursive
calculation incorrectly replaces the projection of a complemented arbitrary
subgame by the negation of its projection. Keep the complemented subgame inside
the projection, or restrict the example to a terminal.

The claim of a unique nontrivial player-symmetric involution is valid if one
requires the same unary Boolean map independently on every coordinate. It is
false for general equivariant involutions: swapping only $0^K$ and $1^K$ while
fixing all other vectors is another such involution for $K\ge2$.

\subsection{Terminal syntax is not a safe semantic quotient}

The cheap-talk quotient in \texttt{solutions.tex:220} need not preserve outcomes.
Let $P_0c.P_1y.\langle y=0,\top\rangle$ and fix $1$'s policy $y:=c$. Changing only
$c$ changes the outcome, although $c$ occurs in no terminal goal. A quotient
needs an explicitly proven observational congruence.

The slicing lemma in \texttt{solutions.tex:495} is false because it switches a
whole role's strategy from the irrelevance of just one occurrence. Let
$P_1v.P_1u.\langle v=0,\top\rangle$. The final $u$-occurrence is irrelevant, but
replacing player 1's full strategy can change $v$. Repair the conclusion to a
local-policy change, or require all of that role's occurrences to be irrelevant.

The definition of game equivalence by visibility-DAG isomorphism must include
the terminal interpretations and typed labels. Identical graphs with different
terminal goals are not equivalent games. Likewise, consistency of every Boolean
combination of terminal goals does not establish strategic noninterference.

\subsection{Composition needs more than acyclic graph gluing}

The splice in \texttt{algebra.tex:704} specifies an arena but no complete terminal
composition. Two independently produced copies of a shared boundary value can
also disagree: each component may determine a valid singleton play while their
join is empty. A cut interface must identify which side produces each value and
which side consumes it, or require a compatibility proof.

The suffix/prefix edge conditions in that section are reversed relative to the
earlier-to-later orientation. Several clock and megaphone definitions are also
incomplete: the clock written before the game is described as after it; a new
role lacks a terminal component; fixed echo behavior is an extra operational
primitive rather than an ordinary unrestricted role choice. These constructions
need definitions before their algebraic laws can be assessed.

Role merging preserves decentralized strategies only when occurrence-level views
remain unchanged. Recomputing a merged role's memory or pooling its observations
can change the game. The phrase ``one coalition acts as a single role'' must say
which operation is intended.

\subsection{Relational representation and counterfactuals}

Totality only on histories reached under a fixed complete profile, as in
\texttt{ra.tex:101}, is sufficient to run that profile. It is insufficient for
deviations. Strategies must be defined on every legal observation that can arise
against alternative opponents, not just the equilibrium path.

Excluding a hidden column from a local key prevents direct access; it does not
prevent indirect signaling through a visible intermediary. The DQBF countermodel
is also a counterexample to a general no-signaling conclusion from key exclusion.
Retaining all local tables, including off-path rows, is essential for the
game-theoretic reading of relational queries.

\subsection{Public properties and session translations}

Exact identifiability in \texttt{extensions.tex:43} misses partial information.
Over four possible secret values, observing parity can inform a later choice
without identifying the secret. Use observation partitions/maps rather than an
all-or-nothing visibility edge. A claimed public property also needs evidence
or an explicit trust/ideal-functionality assumption. An inconsistent public
theory identifies everything. Unique existence alone does not introduce an
$\iota$-term into the first-order signature. With all certificates trivial,
the proposed private-value extension does not recover the broadcast base, so its
stated conservativity claim is not established.

The session translations are at most structural encodings without further
conditions. Ground message sorts do not force repeated sends to carry the same
value, atomic broadcast knowledge is not supplied by a sequence of ordinary
sends, and a private choice can vanish from the target. Refinements checkable in
the global context need not be checkable in the sender's local context. Demanding
all players' goals as refinements also changes a multi-outcome game into a
joint-admissibility problem. The conjecture at \texttt{session.tex:137} needs a
specified equivalence and a corrected source language before it is a theorem
candidate. It is not evidence for a Curry--Howard correspondence as written.

\section{Finite checks and their limits}

\texttt{research/finite\_audit.py} enumerates total Boolean policies, retaining
off-path behavior and replacing all policies owned by a role during deviations.
It confirms the DQBF, hidden-matching, cooperative-negation, cheap-talk, and
slicing examples above. It also checks all 256 payoff tables for two Boolean
moves in each of two observation regimes: 512 games and 3,072 profiles, with 512
broadcast projection comparisons. All checks passed.

These are finite experiments, not Lean/Rocq proofs. The manuscript's
\texttt{interpreter.v} is pseudocode with an unfinished assignment definition and
does not mechanize its theorems. A useful first formalization would implement
the exact game object, role transport, and the finite dependent interaction
theorem; the manuscript's full claim set should not be encoded as axioms.

\end{document}
